% O GIQ, como o resto da Elsevier, quer o abstract em BLOCO UNICO, sem quebra de
% paragrafo, e em ate 250 palavras (contadas com as macros expandidas por
% code/analysis/check_highlights.py). A abertura pela pergunta do gestor e
% deliberada: duas lentes da auditoria externa leram a versao que abria pelo
% protocolo como escrita para um publico de ML.
\begin{abstract}

A municipal environmental manager asks a language model what the binding
PM\textsubscript{2.5} standard is in their country, and nothing in the answer
says whether it matches the official register. We propose a protocol for that
check and report what it finds. Answers are scored against official registers,
by code wherever the correct answer is a published value and by a three-vendor
judge panel elsewhere, across 25 countries, 14 deployment stacks and four
languages (\ncells{} scored cells). The three fixes an agency can deploy without
building anything all fail. Asking in the country's own language lowers accuracy
by \nnativa{} percentage points (mixed model with a random intercept per country,
$p=\nhtwomixedcountryp$) and makes the model \nverifdiffabs{} points less likely
to return a register-checkable value; declaring the official role changes
nothing ($p=\npersonap$); and a regionally fine-tuned 7B model ranks last of
fourteen. Accuracy is lowest on the task with administrative consequence,
recalling the binding value (\ntasktoneshort{} against \nfloorrestshort{} for
synthesis and recommendation). A North/South gap of $\ntiergap$~pp persists
(permutation over countries, $p=\ntiergapperm$), is larger on the outcomes no
judge touches, and is concentrated on tasks whose answer is a national fact; on
the national standard the odds ratio is \norrawone{} with a country-level
interval of $\norrawoneci$. A corpus-coverage mechanism is not established once
inference is carried at the country level. The protocol transfers to any
regulated domain whose correct answer is published.

\end{abstract}
